\section{Introduction}
\label{sec:introduction}

Decision Trees (DTs) and their ensemble extensions, such as Random Forests (RF) \cite{breiman2001random} and Gradient Boosted Decision Trees (GBDT) \cite{ke2017lightgbm,chen2016xgboost}, remain foundational workhorses in competitive machine learning, tabular data modeling, and mission-critical decision support systems. Despite their widespread adoption and desirable properties---including interpretability, robustness to feature scaling, and native handling of heterogeneous data types---training tree-based models on massive, high-dimensional datasets poses severe computational and memory bottlenecks.

The core computational bottleneck in decision tree construction lies in the recursive node splitting procedure. At each internal node containing $N$ training instances and $D$ continuous features, standard greedy algorithms (e.g., CART \cite{breiman1984classification}) sort instance values along all feature dimensions to identify the optimal split $(j^*, \theta^*)$ that maximizes an impurity reduction metric (such as Gini Gain or Information Gain). This exact search requires $\mathcal{O}(D \cdot N \log N)$ time per node. While histogram-based binning approximations (e.g., in LightGBM and FastForest) discretize features into $B$ discrete bins and compute splits in $\mathcal{O}(D \cdot (N + B))$ time, evaluating all $N$ instances across all $K = D \times B$ candidate split points remains prohibitive when scaling to millions of rows and deep tree hierarchies.

\subsection{Multi-Armed Bandit Splitting: Opportunities and Shortcomings}
To break the linear dependence on sample size $N$ during split evaluation, recent research has framed candidate split evaluation as a Multi-Armed Bandit (MAB) Best Arm Identification (BAI) problem \cite{tiwari2022mabsplit,maron1994hoeffding}. In this formulation, each candidate threshold represents an ``arm'', and evaluating the impurity reduction on a random subset of samples represents a stochastic reward pull. By repeatedly drawing minibatches of size $M \ll N$ and tracking statistical confidence bounds, non-competitive candidate splits can be progressively eliminated before processing all training instances at the node.

However, existing MAB-based tree splitting approaches suffer from four fundamental theoretical and practical limitations:
\begin{enumerate}
    \item \textbf{Violation of the i.i.d. Sampling Assumption:} Standard approaches employ classical Hoeffding or empirical Bernstein bounds that presume infinite-population i.i.d. sampling with replacement. In decision tree training, samples are partitioned without replacement across child nodes. Applying i.i.d. bounds under without-replacement subsampling introduces loose confidence intervals and artificial termination delays, failing to collapse uncertainty to zero even when all $N$ available node samples are observed.
    \item \textbf{Monolithic Feature Resolution Overhead:} Real-world datasets often possess hundreds of continuous features, producing tens of thousands of candidate arms per node ($K = D \cdot B$). Evaluating all arms at full resolution ($B = 256$) in parallel imposes massive CPU cache churn and unnecessary statistical exploration on blatantly sub-optimal features.
    \item \textbf{Rigid Minibatch Sizing:} Conventional implementations fix the minibatch size $M_t$ throughout the elimination rounds. When 95\% of candidate arms have already been discarded and only 2--3 top candidates remain, continuing to draw large batches wastes valuable memory bandwidth.
    \item \textbf{Bandit Bookkeeping Overhead on Leaf Nodes:} As trees grow deeper, sample counts at internal nodes shrink drastically ($N < 500$). The administrative overhead of maintaining arm state, tracking running variances, and calculating logarithmic confidence radii exceeds the time required to perform an exact vectorized sort.
\end{enumerate}

\subsection{Our Contributions}
To address these limitations, we introduce \textbf{MABSplit}, a comprehensive, mathematically rigorous, and high-performance framework for decision tree and random forest acceleration. The primary contributions of this work are summarized as follows:

\begin{itemize}
    \item \textbf{Finite-Population Serfling Empirical Bernstein Bounds:} We derive and integrate Bardenet-Maillard and Serfling finite-population concentration bounds for node impurity variance. This guarantees that empirical uncertainty shrinks proportionally to the unobserved finite population fraction $(1 - (n-1)/N)$, rigorously collapsing to zero at full sample exhaustion without PAC-optimality violations.
    \item \textbf{Hierarchical Coarse-to-Fine Arm Pruning (H-MAB):} We introduce a two-tier screening mechanism that first filters uninformative features at coarse 16-bin quantization with a strict Bonferroni budget $\delta_{\text{coarse}} = \delta/3$, subsequently refining candidate arms to 256-bin precision with $\delta_{\text{fine}} = 2\delta/3$, provably preserving the global $(\epsilon, \delta)$-PAC guarantee.
    \item \textbf{Adaptive Dynamic Minibatch Allocation (ADMA):} We develop a sample allocation schedule that scales minibatch size $M_t$ dynamically as a function of the active arm cardinality $|\mathcal{A}_t|$, saving up to 30\% of total sample queries during the terminal phase of arm selection.
    \item \textbf{Depth-Adaptive Hybrid Solver:} We formulate an automatic dispatch mechanism governed by a critical node size threshold $N_{\text{thresh}}$, seamlessly delegating shallow sub-trees to SIMD-vectorized exact sorting and large nodes to the accelerated MAB engine.
    \item \textbf{Split Fidelity Metric and Empirical Validation:} We introduce the Split Agreement Rate metric and validate MABSplit across synthetic and real-world benchmarks (Adult, Covertype, HIGGS), achieving $3.2\times$--$11.8\times$ wall-clock speedups over exact CART with $>99.2\%$ split fidelity.
\end{itemize}
